
\documentclass{article}
\usepackage{colortbl}
\usepackage{makecell}
\usepackage{multirow}
\usepackage{supertabular}

\begin{document}

\newcounter{utterance}

\centering \large Interaction Transcript for game `referencegame', experiment `line\_grids\_columns', episode 3 with reka{-}flash{-}3.1{-}t0.0.
\vspace{24pt}

{ \footnotesize  \setcounter{utterance}{1}
\setlength{\tabcolsep}{0pt}
\begin{supertabular}{c@{$\;$}|p{.15\linewidth}@{}p{.15\linewidth}p{.15\linewidth}p{.15\linewidth}p{.15\linewidth}p{.15\linewidth}}
    \# & $\;$A & \multicolumn{4}{c}{Game Master} & $\;\:$B\\
    \hline

    \theutterance \stepcounter{utterance}  
    & & \multicolumn{4}{p{0.6\linewidth}}{
        \cellcolor[rgb]{0.9,0.9,0.9}{
            \makecell[{{p{\linewidth}}}]{
                \texttt{\tiny{[P1$\langle$GM]}}
                \texttt{You are given three grids, where each of them is 5 by 5 in size.} \\
\texttt{Grids have empty cells marked with "▢" and filled cells marked with "X".} \\
\texttt{Your task is to generate a referring expression that best describes the target grid while distinguishing it from the two other distractor grids.} \\
\texttt{The first grid is the target grid, and the following two grids are the distractors.} \\
\\ 
\texttt{Target grid:} \\
\\ 
\texttt{X □ □ X X} \\
\texttt{X □ □ X X} \\
\texttt{X □ □ X X} \\
\texttt{X □ □ X X} \\
\texttt{X □ □ X X} \\
\\ 
\texttt{Distractor grid 1:} \\
\\ 
\texttt{□ □ X X X} \\
\texttt{□ □ X X X} \\
\texttt{□ □ X X X} \\
\texttt{□ □ X X X} \\
\texttt{□ □ X X X} \\
\\ 
\texttt{Distractor grid 2:} \\
\\ 
\texttt{X □ X X □} \\
\texttt{X □ X X □} \\
\texttt{X □ X X □} \\
\texttt{X □ X X □} \\
\texttt{X □ X X □} \\
\\ 
\texttt{Instruction: Describe the target grid.} \\
\texttt{Generate the referring expression starting with the tag "Expression: " for the given target grid. Omit any other text.} \\
            }
        }
    }
    & \\ \\

    \theutterance \stepcounter{utterance}  
    & \multicolumn{4}{p{0.6\linewidth}}{
        \cellcolor[rgb]{0.9,0.9,0.9}{
            \makecell[{{p{\linewidth}}}]{
                \texttt{\tiny{[P1$\rangle$GM]}}
                \texttt{<reasoning>} \\
\texttt{Let's tackle this problem. I need to create a referring expression that describes the target grid and distinguishes it from the two distractor grids. I'll start by carefully looking at all three grids to understand their differences.} \\
\\ 
\texttt{First, the target grid. Each row has X in the first and last positions, and □ in the middle three. So, it's like a border of X's on the left and right sides. All five rows follow this pattern. So, the structure is two X's at the ends of each row, with empty cells in between.} \\
\\ 
\texttt{Now, Distractor grid 1. Each row here has X's in the last three positions and □ in the first two. So, the X's are on the right side, forming a right{-}aligned block. Each row starts with two empty cells and ends with three X's. That's the opposite of the target's left and right X's.} \\
\\ 
\texttt{Distractor grid 2. Each row here has X's in the first and fourth positions, and □ in the second, third, and fifth. So, it's like a diagonal or a pattern where every other cell is X, but not exactly. Wait, looking at each row: X, □, X, X, □. So, the first and fourth positions are X, but the third is also X. Wait, let me check again. The second row is X □ X X □. So, positions 1, 3, 4 are X? Wait, no. Let me parse each row:} \\
\\ 
\texttt{Row 1: X □ X X □ → positions 1, 3, 4 are X? Wait, no. Let's count. Each row is five cells. For Distractor grid 2:} \\
\\ 
\texttt{Row 1: X (1), □ (2), X (3), X (4), □ (5). So positions 1,3,4 are X. Similarly for other rows. So each row has X in the first, third, and fourth columns. So it's a diagonal pattern but not exactly. It's more like a vertical line on the left and some in the middle. Hmm.} \\
\\ 
\texttt{So the target grid has two X's at the start and end of each row. Distractor 1 has three X's at the end. Distractor 2 has X's in positions 1,3,4} \\
            }
        }
    }
    & & \\ \\

    \theutterance \stepcounter{utterance}  
    & & & \multicolumn{2}{p{0.3\linewidth}}{
        \cellcolor[rgb]{0.9,0.9,0.9}{
            \makecell[{{p{\linewidth}}}]{
                \texttt{\tiny{[GM$|$GM]}}
                \texttt{Invalid generated expression} \\
            }
        }
    }
    & & \\ \\

\end{supertabular}
}

\end{document}
